\chapter{Results and Discussion}

The Sky ResQ Dashboard was tested extensively both in rigorous simulated environments (utilizing the ArduPilot Software-In-The-Loop (SITL) architectures locally via Cygwin) and via physical test flights utilizing long-range 915MHz Silicon Labs SiK telemetry radio sets communicating natively with PX4/Pixhawk flight hardware. The empirical data, latency outcomes, and visual interface deployments are analyzed in the ensuing sections.

\section{User Interface Analysis and Cognitive Reduction}
The culminating design of the primary Ground Control Station interface achieves its overarching goal of drastically minimizing operator cognitive load during mission-critical events. Figure \ref{fig:dashboard_main} presents the primary operational dashboard view, utilized actively during active flight sequences. Fundamentally rejecting the heavily nested configuration tree models indicative of legacy systems (such as Mission Planner), the Sky ResQ HMI employs a flat, grid-aligned, high-contrast modular card geography.

% TODO: Placeholder for Main Dashboard Screenshot
\begin{figure}[htp]
    \centering
    \includegraphics[width=\textwidth,keepaspectratio]{img/main_dashboard.png}
    \caption{Primary Dashboard Interface. The modular layout cleanly segregates Geospatial positioning (left) from Critical Flight Telemetrics (right) minimizing operator saccades.}
    \label{fig:dashboard_main}
\end{figure}

The expansive leftmost geographic pane seamlessly embeds a persistent React-Leaflet \texttt{DroneMap} component, dynamically animating and anchoring focus directly atop the UAV's global coordinates (parsed explicitly from the \texttt{GLOBAL\_POSITION\_INT} `int32` MAVLink indices). The auxiliary right-hand pane clusters volatile data logically: Battery Electrochemical Status, Satellite Positioning Matrices (HDOP, Fix Types), combined VFR vectors (airspeed, ascent vectors, throttle gradients), and the crowning real-time Attitude/Compass visual indicator sets.

\section{Telemetry Integrity, Throughput, and Latency Matrices}
Empirical benchmarks logged during physical hardware connection sequences proved that the custom-tailored TypeScript MAVLink parser algorithm successfully streams and accurately decodes data vectors at 57,600 baud speeds, successfully managing an ingestion and rendering rate averaging 15 to 20 raw MAVLink packets per second. Integration of the Zustand atomic state update architecture structurally negated UI repaint stuttering natively inherent to high-volume bursts within React's standard framework.

Figure \ref{fig:telemetry_latency} mathematically maps the sequential flow of JSON data transitioning across the strictly typed Electron IPC (Inter-Process Communication) boundaries. Against traditional networking configurations like local WebSocket topologies or UDP sockets, the deployed direct-IPC channels presented a near non-existent transmission bottleneck, reliably tracking less than $<1.5$ms variance between the dedicated Node.js background \texttt{CP2102} buffer ingress thread and the Chromium visual rendering cascade.

% TODO: Placeholder for Latency/Performance Graph or IPC Flow Diagram
\begin{figure}[htp]
    \centering
    \includegraphics[width=\textwidth,keepaspectratio]{img/ipc.png}
    \caption{Data Transfer Latency tracking measurements traversing the Node.js/Chromium Electron IPC Boundary versus conventional REST models.}
    \label{fig:telemetry_latency}
\end{figure}

\section{Cross-Platform Native Desktop Viability and 'Kiosk' Operation}
Embracing the overarching Electron compilation configuration, the source code natively synthesized stand-alone application executables demanding zero environmental prerequisites for deployment. The pipeline consistently outputted fully independent distribution binaries specifically formatted for Windows environments (\texttt{.exe}), Linux builds (\texttt{.AppImage}/deb derivatives), and macOS packages (\texttt{.dmg}). This achievement absolutely confirmed the project’s core objective regarding total hardware architecture and Operating System independence.

Showcased in Figure \ref{fig:electron_window}, the environment physically detaches from the restricting confines of standard tabbed browser operation, launching as an isolated native application class. This structural capability supports "kiosk-mode" scenarios, strictly segregating the GCS application scope from potentially disruptive arbitrary host OS notifications or background tab processing, a critical capability when utilizing the system on ruggedized off-grid deployment tablets.

% TODO: Placeholder for Electron Native App Screenshot
\begin{figure}[htp]
    \centering
    \includegraphics[width=\textwidth,keepaspectratio]{img/electron_native.png}
    \caption{Sky ResQ Dashboard deployed locally as a fully integrated native OS executable via Electron}
    \label{fig:electron_window}
\end{figure}

\section{Algorithm Robustness: Recovery from Structural Hardware Disruptions}
Ultimately, the final architectural revisions systematically eliminated the acute limitations identified in generic HTML5 Web Serial API implementations, notably concerning long-term, high-voltage continuous stream survivability. 

The custom-engineered exponential backoff and buffer-overrun interception subroutines empirically validated their necessity during localized USB cable instability trials. In scenarios engineered to mimic momentary physical vibrations, aggressive EMI interference, or OS-level Thread Suspensions (i.e. forcibly collapsing a laptop lid), the GCS demonstrated a self-healing paradigm. The algorithm accurately detects the stream disruption, forcefully destroys the dangling readable lock indices, flushes the localized queue sizes, and automatically re-initiates the connection socket seamlessly—restoring optimal telemetry flow entirely transparent to the operator, affirming the Sky ResQ Ground Control Station as an objectively robust tool suite adaptable for austere environments.
